% Options for packages loaded elsewhere
\PassOptionsToPackage{unicode}{hyperref}
\PassOptionsToPackage{hyphens}{url}
%
\documentclass[
  11pt,
  ignorenonframetext,
]{beamer}
\usepackage{pgfpages}
\setbeamertemplate{caption}[numbered]
\setbeamertemplate{caption label separator}{: }
\setbeamercolor{caption name}{fg=normal text.fg}
\beamertemplatenavigationsymbolsempty
% Prevent slide breaks in the middle of a paragraph
\widowpenalties 1 10000
\raggedbottom
\setbeamertemplate{part page}{
  \centering
  \begin{beamercolorbox}[sep=16pt,center]{part title}
    \usebeamerfont{part title}\insertpart\par
  \end{beamercolorbox}
}
\setbeamertemplate{section page}{
  \centering
  \begin{beamercolorbox}[sep=12pt,center]{part title}
    \usebeamerfont{section title}\insertsection\par
  \end{beamercolorbox}
}
\setbeamertemplate{subsection page}{
  \centering
  \begin{beamercolorbox}[sep=8pt,center]{part title}
    \usebeamerfont{subsection title}\insertsubsection\par
  \end{beamercolorbox}
}
\AtBeginPart{
  \frame{\partpage}
}
\AtBeginSection{
  \ifbibliography
  \else
    \frame{\sectionpage}
  \fi
}
\AtBeginSubsection{
  \frame{\subsectionpage}
}
\usepackage{amsmath,amssymb}
\usepackage{iftex}
\ifPDFTeX
  \usepackage[T1]{fontenc}
  \usepackage[utf8]{inputenc}
  \usepackage{textcomp} % provide euro and other symbols
\else % if luatex or xetex
  \usepackage{unicode-math} % this also loads fontspec
  \defaultfontfeatures{Scale=MatchLowercase}
  \defaultfontfeatures[\rmfamily]{Ligatures=TeX,Scale=1}
\fi
\usepackage{lmodern}
\usetheme[]{boxes}
\ifPDFTeX\else
  % xetex/luatex font selection
\fi
% Use upquote if available, for straight quotes in verbatim environments
\IfFileExists{upquote.sty}{\usepackage{upquote}}{}
\IfFileExists{microtype.sty}{% use microtype if available
  \usepackage[]{microtype}
  \UseMicrotypeSet[protrusion]{basicmath} % disable protrusion for tt fonts
}{}
\makeatletter
\@ifundefined{KOMAClassName}{% if non-KOMA class
  \IfFileExists{parskip.sty}{%
    \usepackage{parskip}
  }{% else
    \setlength{\parindent}{0pt}
    \setlength{\parskip}{6pt plus 2pt minus 1pt}}
}{% if KOMA class
  \KOMAoptions{parskip=half}}
\makeatother
\usepackage{xcolor}
\newif\ifbibliography
\setlength{\emergencystretch}{3em} % prevent overfull lines
\providecommand{\tightlist}{%
  \setlength{\itemsep}{0pt}\setlength{\parskip}{0pt}}
\setcounter{secnumdepth}{-\maxdimen} % remove section numbering
\addtobeamertemplate{title page}{\centering \includegraphics[width=1\textwidth]{images/banner12.png}}{}
\ifLuaTeX
  \usepackage{selnolig}  % disable illegal ligatures
\fi
\usepackage{bookmark}
\IfFileExists{xurl.sty}{\usepackage{xurl}}{} % add URL line breaks if available
\urlstyle{same}
\hypersetup{
  pdftitle={Senses \& Cognition},
  hidelinks,
  pdfcreator={LaTeX via pandoc}}

\title{Senses \& Cognition}
\author{}
\date{\vspace{-2.5em}Office: 305 Lewis Hall\\
Dr.~Ethan Linck}

\begin{document}
\frame{\titlepage}

\begin{frame}{Birds are visual organisms}
\phantomsection\label{birds-are-visual-organisms}
\begin{itemize}
\tightlist
\item
  Vision deeply important to feeding, reproduction, avoiding predators
\item
  Bird eyes are large (human-sized in some eagles and owls) and
  sophisticated
\item
  Songbirds and raptors can resolve details 2-3x farther than us
\end{itemize}

\begin{figure}
\includegraphics[width = 0.6\linewidth]{images/bluebop.png}
\end{figure}

\hfill \tiny \href{https://macaulaylibrary.org/asset/608319877}{\emph{Paradisornis
rudolphi}}
\end{frame}

\begin{frame}{Bird eye anatomy}
\phantomsection\label{bird-eye-anatomy}
\begin{itemize}
\tightlist
\item
  both \textbf{lens} and \textbf{cornea} change shape to focus
\item
  focus via \textbf{Campton's} \& \textbf{Brucke's} muscles on
  \textbf{sclerotic ring}
\item
  \textbf{pecten} structure enhances blood supply without constraining
  receptor density
\item
  \textbf{fovea} a pit of high cone density for sharp central vision
\end{itemize}

\begin{figure}
\includegraphics[width = 0.55\linewidth]{images/bird_eye_anatomy.jpg}
\end{figure}

\hfill \tiny \href{https://en.wikipedia.org/wiki/Pecten_oculi}{Wikipedia}
\end{frame}

\begin{frame}{External eye features}
\phantomsection\label{external-eye-features}
\begin{itemize}
\tightlist
\item
  \textbf{nictating membrane} (L): shared with other reptiles;
  transparent membrane that moves front to back to moisten and protext
  eyes (particularly diving birds)
\item
  \textbf{tapetum lucidum} (R): ``eyeshine'' of night birds caused by
  lipid spheres in retina that scatter light to improve vision;
  physiological, not permanent
\end{itemize}

\begin{figure}
\includegraphics[width = \linewidth]{images/eye_stuff.png}
\end{figure}

\hfill \tiny \href{https://www.audubon.org/news/birds-have-built-goggles}{Audubon};
\href{https://news.maryland.gov/dnr/2024/10/01/native-animal-profile-eastern-whip-poor-will/}{Maryland
Biodiversity Project}
\end{frame}

\begin{frame}{Eye position impacts field of vision}
\phantomsection\label{eye-position-impacts-field-of-vision}
\begin{itemize}
\tightlist
\item
  \textbf{binocular vision}: better depth perceptinon
\item
  \textbf{monocular vision}: can enhance field of view; head bobbing or
  switching angles to compensate for reduced depth perception
\end{itemize}

\begin{figure}
\includegraphics[width = 0.6\linewidth]{images/vision.png}
\end{figure}

\hfill \tiny \href{https://macaulaylibrary.org/asset/61674401}{\emph{Columba
livia}}; \href{https://macaulaylibrary.org/asset/63942211}{\emph{Bubo
virginianus}};
\href{https://en.wikipedia.org/wiki/Bird_vision}{Wikipedia}
\end{frame}

\begin{frame}{Eye position impacts field of vision}
\phantomsection\label{eye-position-impacts-field-of-vision-1}
\begin{figure}
\includegraphics[width = 0.8\linewidth]{images/field_of_vision.jpg}
\end{figure}

\hfill \tiny (Martin 2022)
\end{frame}

\begin{frame}{Light sensors}
\phantomsection\label{light-sensors}
\begin{itemize}
\tightlist
\item
  \textbf{rods} are highly sensitive receptors for low light
\item
  \textbf{cones} are coiled proteins containing carotenoid pigment;
  sensitive to different portions of the visible spectrum
\item
  \textbf{oil droples} reduce light but enhance color discrimination
\end{itemize}

\begin{figure}
\includegraphics[width = 0.9\linewidth]{images/cones2.png}
\end{figure}

\hfill \tiny (Toomey et al.~2016)
\end{frame}

\begin{frame}{Birds perceive ultraviolet light}
\phantomsection\label{birds-perceive-ultraviolet-light}
\begin{itemize}
\tightlist
\item
  Birds can see light from 325 to 700 nanometers (versus 400 to 700 in
  humans)
\item
  Cone types for red, green, blue, and ultraviolet light (versus just
  RGB in humans)
\end{itemize}

\begin{figure}
\includegraphics[width = 0.7\linewidth]{images/bird_vision_spectra.png}
\end{figure}

\hfill \tiny (\href{Bird_vision}{Wikipedia})
\end{frame}

\begin{frame}{Tetrahedral color space}
\phantomsection\label{tetrahedral-color-space}
\begin{itemize}
\tightlist
\item
  Combinations of cones create 2D color space in humans (b), 3D color
  space in birds (a)
\item
  Color comprised of \textbf{hue} (angle from center of color vector)
  and saturation (distance from center)
\item
  Vertical axis of avian color tetrahedron defines colors that are
  unperceivable and unimaginable to humans
\end{itemize}

\begin{figure}
\includegraphics[width = 0.7\linewidth]{images/tetra.jpg}
\end{figure}

\hfill \tiny (Janas et al.~2019)
\end{frame}

\begin{frame}{UV colors impact reproduction}
\phantomsection\label{uv-colors-impact-reproduction}
\begin{itemize}
\tightlist
\item
  Male Bluethroats (\emph{Luscinia svecica}) were left untreated (top),
  rubbed with baby oil (middle), or rubbed with sunblock (bottom)
\item
  Birds treated with sunblock had lower mating success
\end{itemize}

\begin{figure}
\includegraphics[width = 0.8\linewidth]{images/bluethroat.png}
\end{figure}

\hfill \tiny (\href{https://macaulaylibrary.org/asset/182455611}{ML182455611};
Johnsen et al.~1998)
\end{frame}

\begin{frame}{Bird hearing is not extraordinary}
\phantomsection\label{bird-hearing-is-not-extraordinary}
\begin{figure}
\includegraphics[width = 0.7\linewidth]{images/hearing.png}
\end{figure}

\hfill \tiny (Gill et al.~2019)
\end{frame}

\begin{frame}{Bird ear anatomy}
\phantomsection\label{bird-ear-anatomy}
\begin{itemize}
\tightlist
\item
  no external ear (\textbf{pinnae} of mammals); instead, an \textbf{ear
  funnel} covered with feathers
\item
  the \textbf{stapes} (\textbf{columella}) transmits sound vibrations to
  the \textbf{tympanum}
\item
  the \textbf{semicircular canals} control balance
\item
  the \textbf{cochlea} is the inner ear
\end{itemize}

\begin{figure}
\includegraphics[width = 0.4\linewidth]{images/stapes2.png}
\end{figure}

\hfill \tiny (Gill et al.~2019)
\end{frame}

\begin{frame}{Asymmetrical ears aid prey location}
\phantomsection\label{asymmetrical-ears-aid-prey-location}
\begin{figure}
\includegraphics[width = 0.9\linewidth]{images/owl_asymmetry.jpg}
\end{figure}

\hfill \tiny (\emph{Aegolius funereus},
\href{https://www.internationalowlcenter.org/}{International Owl
Center})
\end{frame}

\begin{frame}{Echolocation}
\phantomsection\label{echolocation}
\begin{itemize}
\tightlist
\item
  known in swiftlets (\emph{Aerodromus} and \emph{Collocalia}) and
  oilbirds (\emph{Steatornis})
\end{itemize}

\begin{figure}
\includegraphics[width = 0.7\linewidth]{images/steatornis.png}
\end{figure}

\url{https://www.youtube.com/watch?v=lKEYMJ1vjj8}

\hfill \tiny \href{https://macaulaylibrary.org/asset/608319877}{\emph{Steatornis
caripensis}}
\end{frame}

\begin{frame}{Tactile sensation}
\phantomsection\label{tactile-sensation}
\begin{itemize}
\tightlist
\item
  \textbf{Hebst corpuscles}: largest and most common touch receptors
\item
  Merkel corpuscles: skin, bill, tongue
\item
  Gandry corpuscles: in duck bills only
\end{itemize}

\begin{figure}
\includegraphics[width = \linewidth]{images/hebst.png}
\end{figure}

\hfill \tiny (Crole \& Soley 2014)
\end{frame}

\begin{frame}{Birds have a poorly developed sense of taste}
\phantomsection\label{birds-have-a-poorly-developed-sense-of-taste}
\begin{itemize}
\tightlist
\item
  2 (chicken) to 62 (Japanese Quail) taste buds on rear of tongue and
  back of throat versys 10,000+ on tongue in humans
\end{itemize}

\begin{figure}
\includegraphics[width = \linewidth]{images/tastebuds.jpg}
\end{figure}

\hfill \tiny (Crole \& Soley 2014)
\end{frame}

\begin{frame}{Birds are not sensitive to capsaicin}
\phantomsection\label{birds-are-not-sensitive-to-capsaicin}
\begin{itemize}
\tightlist
\item
  \textbf{directed deterrence hypothesis}: secondary metabolites in ripe
  fruit function to deter fruit consumption by vertebrates that do not
  disperse seeds, while not impacting consumption by those that do
\item
  Levey et al.: NO chili consumption by rodents in AZ / Bolivia, even
  though shown in lab; entirely by birds
\end{itemize}

\begin{figure}
\includegraphics[width = 0.5\linewidth]{images/chilis.png}
\end{figure}

\hfill \tiny (Levey et al.~2006)
\end{frame}

\begin{frame}{Hummingbirds can taste sugar}
\phantomsection\label{hummingbirds-can-taste-sugar}
\begin{itemize}
\tightlist
\item
  strong feeding preference for sucrose and erythritol, which activate
  T1R1-T1R3 sensors
\item
  limited time feeding on water or aspartame, which fail to activate
  T1R1-T1R3 sensors
\end{itemize}

\begin{figure}
\includegraphics[width = 0.8\linewidth]{images/sweet_hummer.png}
\end{figure}

\hfill \tiny (Baldwin et al.~2014)
\end{frame}

\begin{frame}{Sweet receptor a novel modification of umami T1R3 protein}
\phantomsection\label{sweet-receptor-a-novel-modification-of-umami-t1r3-protein}
\begin{itemize}
\tightlist
\item
  Most birds lack T1R2 sweet receptor
\item
  Hummingbirds have modified TR13 protein that responds to carbohydrates
\end{itemize}

\begin{figure}
\includegraphics[width = 0.8\linewidth]{images/t1r.png}
\end{figure}

\hfill \tiny (Baldwin et al.~2014)
\end{frame}

\begin{frame}{Enhanced olfactory bulb in Kiwi}
\phantomsection\label{enhanced-olfactory-bulb-in-kiwi}
\begin{itemize}
\tightlist
\item
  Olfactory bulb large in scent-dependent species: kiwis, vultures,
  petrels
\end{itemize}

\begin{figure}
\includegraphics[width = 0.9\linewidth]{images/olfactory.png}
\end{figure}

\hfill \tiny (Baldwin et al.~2014)
\end{frame}

\begin{frame}{Birds can be very intelligent}
\phantomsection\label{birds-can-be-very-intelligent}
\begin{itemize}
\tightlist
\item
  Brain large relative to body size, particularly in songbirds and
  parrots
\item
  The smartest species show object permanence, can delay gratification,
  remember past episodes and plan ahead, reason to solve puzzles,
  recognize themselves in the mirror, infer mental states of others,
  show awareness of own learning (metacognition), learn complex songs
\end{itemize}

\begin{figure}
\includegraphics[width = \linewidth]{images/brain_scaling.png}
\end{figure}

\hfill \tiny (Olkowicz et al.~2016)
\end{frame}

\begin{frame}{Bird brain anatomy}
\phantomsection\label{bird-brain-anatomy}
\begin{itemize}
\tightlist
\item
  \textbf{forebrain}: instinctual behaviors, learning, sensory
  integration
\item
  \textbf{midbrain}: vision, coordination, physiology, hormones
\item
  \textbf{hindbrain / medulla}: links to spinal cord and peripheral
  nervous system
\end{itemize}

\begin{figure}
\includegraphics[width = 0.5\linewidth]{images/brain_anatomy.png}
\end{figure}

\hfill \tiny (Gill et al.~2019)
\end{frame}

\begin{frame}{An evolving view of the avian forebrain}
\phantomsection\label{an-evolving-view-of-the-avian-forebrain}
\begin{itemize}
\tightlist
\item
  Majority of avian forebrain long thought to be homologous with
  mammalian \textbf{striatum} (involved in instinctive behavior)
\item
  More recently discovered to consist of \textbf{pallial} elements
  associated with learning and intelligence
\end{itemize}

\begin{figure}
\includegraphics[width = 0.8\linewidth]{images/brain_evol_theory.png}
\end{figure}

\hfill \tiny (Emery 2007)
\end{frame}

\begin{frame}{Birds have high neuronal density}
\phantomsection\label{birds-have-high-neuronal-density}
\begin{itemize}
\tightlist
\item
  Corvids and parrots have more pallial neurons in forebrains than
  primates 3x as large!
\end{itemize}

\begin{figure}
\includegraphics[width = \linewidth]{images/neuron_density.png}
\end{figure}

\hfill \tiny (Olkowicz et al.~2016)
\end{frame}

\begin{frame}{Seed caching in Clark's Nutcrackers}
\phantomsection\label{seed-caching-in-clarks-nutcrackers}
\begin{itemize}
\tightlist
\item
  Seed caching behavior requires impressive cognitive ability
  (especially spatial memory)
\item
  Experiments show Clark's Nutcrackers recover stashed seeds at a higher
  rate than random expectations (but this decays over time)
\end{itemize}

\begin{figure}
\includegraphics[width = 0.9\linewidth]{images/memory.png}
\end{figure}

\hfill \tiny \href{https://macaulaylibrary.org/asset/51262821}{\emph{Nucifraga
columbiana}}; Balda \& Kamil 1992
\end{frame}

\begin{frame}{New Caledonian crow tool use}
\phantomsection\label{new-caledonian-crow-tool-use}
\begin{figure}
\includegraphics[width = 0.6\linewidth]{images/crows.png}
\end{figure}

\url{https://www.youtube.com/watch?v=EryZPmOxwC0&t=15s}

\hfill \tiny (Bayern et al.~2018;
\href{https://www.minervateam.hu/contactl/crow-slot/crows-make-complex-tools-6882266/}{Minerva})
\end{frame}

\begin{frame}{Keas understand probability}
\phantomsection\label{keas-understand-probability}
\begin{figure}
\includegraphics[width = 0.9\linewidth]{images/kea.png}
\end{figure}

\url{https://www.youtube.com/watch?v=Yj718A7_s4A}

\hfill \tiny \href{https://macaulaylibrary.org/asset/57710841}{\emph{Nestor
notabilis}}
\end{frame}

\begin{frame}{Alex the Parrot}
\phantomsection\label{alex-the-parrot}
\begin{figure}
\includegraphics[width = 0.6\linewidth]{images/alex.png}
\end{figure}

\url{https://www.youtube.com/watch?v=ldYkFdu5FJk}
\end{frame}

\begin{frame}{References}
\phantomsection\label{references}
\begin{itemize}
\item
  Title slide species: \emph{Aviceda subcristata}, \emph{Cacatua
  moluccensis}, \emph{Ardena creatopus}, \emph{Phyllastrephus debilis}
\item
  Martin, G. R. (2022). Avian vision. Current Biology, 32(20),
  R1079-R1085.
\item
  Toomey, M. B., Lind, O., Frederiksen, R., Curley Jr, R. W., Riedl, K.
  M., Wilby, D., \ldots{} \& Corbo, J. C. (2016). Complementary shifts
  in photoreceptor spectral tuning unlock the full adaptive potential of
  ultraviolet vision in birds. Elife, 5, e15675.
\item
  Janas, K., Podmokła, E., Lutyk, D., Dubiec, A., Gustafsson, L.,
  Cichoń, M., \& Drobniak, S. (2018). Influence of haemosporidian
  infection status on structural and carotenoid‐based colouration in the
  blue tit Cyanistes caeruleus. Journal of Avian Biology, 49(10),
  e01840.
\end{itemize}
\end{frame}

\begin{frame}{References}
\phantomsection\label{references-1}
\begin{itemize}
\item
  Johnsen, A., Andersson, S., Örnborg, J., \& Lifjeld, J. T. (1998).
  Ultraviolet plumage ornamentation affects social mate choice and sperm
  competition in bluethroats (Aves: Luscinia s. svecica): a field
  experiment. Proceedings of the Royal Society of London. Series B:
  Biological Sciences, 265(1403), 1313-1318.
\item
  Gill, F. B., Prum, R. O., \& Robinson, S. K. (2019). Ornithology (4th
  ed.). W. H. Freeman.
\item
  Tucker, A. S. (2017). Major evolutionary transitions and innovations:
  the tympanic middle ear. Philosophical Transactions of the Royal
  Society B: Biological Sciences, 372(1713), 20150483.
\item
  Crole, M. R., \& Soley, J. T. (2014). Comparative distribution and
  arrangement of Herbst corpuscles in the oropharynx of the ostrich
  (Struthio camelus) and emu (Dromaius novaehollandiae). The Anatomical
  Record, 297(7), 1338-1348.
\item
  Niknafs, S., Navarro, M., Schneider, E. R., \& Roura, E. (2023). The
  avian taste system. Frontiers in Physiology, 14, 1235377.
\end{itemize}
\end{frame}

\begin{frame}{References}
\phantomsection\label{references-2}
\begin{itemize}
\item
  Levey, D. J., Tewksbury, J. J., Cipollini, M. L., \& Carlo, T. A.
  (2006). A field test of the directed deterrence hypothesis in two
  species of wild chili. Oecologia, 150(1), 61-68.
\item
  Baldwin, M. W., Toda, Y., Nakagita, T., O'Connell, M. J., Klasing, K.
  C., Misaka, T., \ldots{} \& Liberles, S. D. (2014). Evolution of sweet
  taste perception in hummingbirds by transformation of the ancestral
  umami receptor. Science, 345(6199), 929-933.
\item
  Corfield, J. R., Eisthen, H. L., Iwaniuk, A. N., \& Parsons, S.
  (2014). Anatomical specializations for enhanced olfactory sensitivity
  in kiwi, Apteryx mantelli. Brain Behavior and Evolution, 84(3),
  214-226.
\item
  Olkowicz, S., Kocourek, M., Lučan, R. K., Porteš, M., Fitch, W. T.,
  Herculano-Houzel, S., \& Němec, P. (2016). Birds have primate-like
  numbers of neurons in the forebrain. Proceedings of the National
  Academy of Sciences, 113(26), 7255-7260.
\end{itemize}
\end{frame}

\begin{frame}{References}
\phantomsection\label{references-3}
\begin{itemize}
\item
  Emery, N. J. (2007). Evolution of the avian brain and intelligence.
  Current Biology, 15, 46.
\item
  Balda, R. P., \& Kamil, A. C. (1992). Long-term spatial memory in
  Clark's nutcracker, Nucifraga columbiana. Animal Behaviour, 44(4),
  761-769.
\item
  Bayern, A. V., Danel, S., Auersperg, A. M. I., Mioduszewska, B., \&
  Kacelnik, A. (2018). Compound tool construction by New Caledonian
  crows. Scientific reports, 8(1), 15676.
\end{itemize}
\end{frame}

\end{document}
